\documentclass[10pt,a4paper]{amsart}
\usepackage[margin=19mm]{geometry}
\usepackage{amsmath,amssymb,mathtools}
\usepackage[scr=rsfs,cal=euler]{mathalfa}
\usepackage{xfrac}
\usepackage[unicode,hidelinks,bookmarks=false]{hyperref}
\newcommand{\ie}{i.e.\ }
\newcommand{\cf}{cf.\ }
\newcommand{\resp}{respectively\ }
\newcommand{\Subsection}{\subsection}
\providecommand{\nobreakdash}{\nobreak\hbox{-}}
\setlength{\emergencystretch}{2em}
\multlinegap0pt
\usepackage{amssymb}
\usepackage{mathtools}
\usepackage[shortlabels]{enumitem}
\usepackage{xcolor}
\newtheorem{lemma}{Lemma}
\newcommand{\C}{\mathbb C}
\newcommand{\qpoch}[2]{(#1;#2)_\infty}
\title{Complex spectrum of the partial theta function}
\author{Boris Shapiro}
\date{Source: arXiv:2605.29991v2 (CC BY 4.0)}
\begin{document}
\maketitle

\noindent\emph{Source and licence.} Complex spectrum of the partial theta function. Version arXiv:2605.29991v2 (CC BY 4.0), distributed by arXiv under the Creative Commons Attribution 4.0 International licence, http://creativecommons.org/licenses/by/4.0/, as recorded in the arXiv licence metadata for this exact version and shown on the abstract page https://arxiv.org/abs/2605.29991v2. Full licence notice: \texttt{licenses/arxiv-2605.29991-CC-BY-4.0.txt}.\par\smallskip

\noindent\textit{Editorial scope.} Counted unit: the article's lemma lem:cusp-asymptotics-final (source0.tex lines 639--674). The article's complete original proof of it is delivered with it (source0.tex lines 676--777, one \texttt{proof} environment of 102 lines). No mathematics is written, repaired or re-derived: the statement and the proof are the article's own lines, in the article's own notation, and no retained line is substituted or re-flowed.  Retained uncounted as the source-local context the proof needs: lines 199--204; lines 630--632.  Nothing was omitted from any retained span, so the package carries no omission record.  No retained line contains a citation, so the wrapper carries no bibliography and no bibliography entry.  No retained line contains a figure environment, an image-inclusion command or drawing code, so no image is delivered and the package declares no asset dependency.  Editorial formatting only, in addition: an AMS \texttt{amsart} preamble that restates the article's own declarations and macros used by the retained lines, namely the environments \texttt{lemma} on separate counters, together with the standard packages those lines need.  The retained environments print in this document as Lemma 1 in order of appearance, whereas the article prints the same items with the numbering of its own full document.  The complete native source of the article as distributed in the arXiv e-print for this exact version is bundled unchanged as \texttt{sources/201645/source0.tex}.
\begin{lemma}[Jacobi product form]\label{lem:eta-cubed}
For \(|q|<1\),
\[
        \psi(q)=\prod_{\nu=1}^{\infty}(1-q^\nu)^3.
\]
\end{lemma}

\begin{equation}\label{eq:H-product-final}
        H(q,u)=\prod_{\nu=1}^{\infty}(1-q^\nu)(1-e^uq^\nu)(1-e^{-u}q^{\nu-1}),
\end{equation}

\begin{lemma}[Uniform sine and eta asymptotics]\label{lem:cusp-asymptotics-final}
Fix $C>0$ and a compact set $V\Subset\C$.  Uniformly for
\[
        |\tau-\tau_0|\le C\frac{\log N}{N},
        \qquad v\in V,
\]
one has
\begin{align}
        H(q_N(\tau),v/N)
        &=\frac{\psi(q_N(\tau))}{N}
          \left(S_b(\tau,v)+O\left(\frac{\log N}{N}\right)\right),\label{eq:H-sine-final}\\
        \partial_uH(q_N(\tau),v/N)
        &=\psi(q_N(\tau))
          \left(\partial_vS_b(\tau,v)+O\left(\frac{\log N}{N}\right)\right),\label{eq:Hu-sine-final}
\end{align}
where
\begin{equation}\label{eq:Sb-final-def}
        S_b(\tau,v):=\frac{b\tau}{\pi}\sin\frac{\pi v}{b\tau}.
\end{equation}
Moreover, after fixing a residue class of $N$ modulo $2b$,
\begin{equation}\label{eq:Lambda-final-asym}
        \Lambda_N(\tau):=\frac{N(-1)^Nq_N(\tau)^{N(N+1)/2}}{\psi(q_N(\tau))}
        =A_N(\tau)\exp\{Nf_b(\tau)\},
\end{equation}
where
\begin{equation}\label{eq:fb-def}
        f_b(\tau)=\frac{\pi^2}{2b^2\tau}-\frac{\tau}{2},
\end{equation}
$A_N$ is holomorphic and non-vanishing in the same window, and
\begin{equation}\label{eq:AN-bound}
        \log |A_N(\tau)|=O(\log N),
        \qquad
        \partial_\tau\log A_N(\tau)=O(\log N)
\end{equation}
there, after a holomorphic choice of the logarithm.  The estimates also hold after one $\tau$- and $v$-derivative in \eqref{eq:H-sine-final}--\eqref{eq:Hu-sine-final}.
\end{lemma}

\begin{proof}
We use the standard notation
\[
        \qpoch{a}{Q}:=\prod_{k=0}^{\infty}(1-aQ^k).
\]
From the product formula \eqref{eq:H-product-final} and Lemma~\ref{lem:eta-cubed},
\begin{equation}\label{eq:qpoch-quotient}
        \frac{H(q,u)}{\psi(q)}=
        \frac{\qpoch{e^u q}{q}\qpoch{e^{-u}}{q}}{\qpoch{q}{q}^2}.
\end{equation}
Put $Q=q^b=e^{-b\tau/N}$ and $h=b\tau/N$.  Grouping the products in \eqref{eq:qpoch-quotient} into residue classes modulo $b$ gives a singular factor
\begin{equation}\label{eq:singular-factor}
        G_N(\tau,v):=
        \frac{\qpoch{e^{v/N}Q}{Q}\qpoch{e^{-v/N}}{Q}}{\qpoch{Q}{Q}^2}
\end{equation}
and a product $U_N(\tau,v)$ over the non-zero residue classes.  For $a=1,\ldots,b-1$, the corresponding factor of $U_N$ is
\[
        \frac{\qpoch{e^{v/N}q^a}{Q}\qpoch{e^{-v/N}q^a}{Q}}{\qpoch{q^a}{Q}^2}.
\]
Since $q^a$ stays uniformly away from $1$ in the present window, the logarithm of this factor is obtained by a Taylor expansion in $v/N$.  The linear terms cancel and the quadratic remainder is bounded by
\[
        O\left(N^{-2}\sum_{k\ge0}e^{-c k/N}\right)=O(N^{-1}),
\]
uniformly for $v\in V$.  Hence
\begin{equation}\label{eq:UN-estimate}
        U_N(\tau,v)=1+O(N^{-1}),
\end{equation}
with the same bound after one $\tau$- or $v$-derivative, after Cauchy estimates in slightly enlarged compact sets.

It remains to evaluate the singular factor.  Write
\[
        \alpha=\frac{v}{b\tau},
        \qquad Q=e^{-h}.
\]
Then $e^{-v/N}=Q^{\alpha}$ and $e^{v/N}Q=Q^{1-\alpha}$, so
\[
        G_N(\tau,v)=
        \frac{\qpoch{Q^\alpha}{Q}\qpoch{Q^{1-\alpha}}{Q}}{\qpoch{Q}{Q}^2}.
\]
Using the $q$-gamma function
\[
        \Gamma_Q(z)=(1-Q)^{1-z}\frac{\qpoch{Q}{Q}}{\qpoch{Q^z}{Q}},
\]
we get
\[
        G_N(\tau,v)=\frac{1-Q}{\Gamma_Q(\alpha)\Gamma_Q(1-\alpha)}.
\]
As $Q\to1$ inside a fixed Stolz angle, $\Gamma_Q(z)\to\Gamma(z)$ locally uniformly in $z$ after taking reciprocals, and therefore, by Euler's reflection formula,
\[
        G_N(\tau,v)=
        (1-Q)\frac{\sin(\pi\alpha)}{\pi}
        +O(N^{-2})
        =\frac1N\frac{b\tau}{\pi}\sin\frac{\pi v}{b\tau}
          +O(N^{-2}),
\]
locally uniformly in $v$ and $\tau$.  Together with \eqref{eq:UN-estimate}, this proves \eqref{eq:H-sine-final}; differentiating with respect to $u$ is the same as applying $N\partial_v$, which gives \eqref{eq:Hu-sine-final}.  The stated derivative estimates follow from the same locally uniform holomorphic convergence and Cauchy estimates.

We next estimate $\psi(q_N(\tau))=\qpoch{q_N(\tau)}{q_N(\tau)}^3$.  We use the standard Euler--Maclaurin expansion, uniform in the same logarithmic window,
\begin{equation}\label{eq:qpoch-dilog-asym}
        \log\qpoch{a e^{-\gamma/N}}{e^{-b\tau/N}}
        =-\frac{N}{b\tau}\operatorname{Li}_2(a)+O(\log N),
\end{equation}
for $a$ in a compact subset of $\C\setminus[1,\infty)$, and the classical singular case
\begin{equation}\label{eq:singular-qpoch-asym}
        \log\qpoch{e^{-b\tau/N}}{e^{-b\tau/N}}
        =-\frac{\pi^2N}{6b\tau}+O(\log N).
\end{equation}
Both estimates follow by applying Euler--Maclaurin to the logarithm of the product; differentiating the same expansion gives the corresponding $\tau$-derivative with an $O(\log N)$ remainder after the leading term is removed.

Now factor
\[
        \qpoch{q}{q}=\qpoch{q^b}{q^b}
        \prod_{a=1}^{b-1}\qpoch{q^a}{q^b}.
\]
Since
\[
        \sum_{a=0}^{b-1}\operatorname{Li}_2(\omega^a)=\frac{\pi^2}{6b},
        \qquad
        \sum_{a=1}^{b-1}\operatorname{Li}_2(\omega^a)=\frac{\pi^2}{6b}-\frac{\pi^2}{6},
\]
\eqref{eq:qpoch-dilog-asym} and \eqref{eq:singular-qpoch-asym} give
\[
        \log\qpoch{q_N(\tau)}{q_N(\tau)}
        =-\frac{\pi^2N}{6b^2\tau}+O(\log N).
\]
Therefore
\begin{equation}\label{eq:psi-cusp-final}
        \log\psi(q_N(\tau))
        =-\frac{\pi^2N}{2b^2\tau}+O(\log N),
\end{equation}
again with the analogous derivative estimate after subtracting the displayed leading term.
Finally,
\[
        q_N(\tau)^{N(N+1)/2}
        =\omega^{N(N+1)/2}\exp\left(-\frac{\tau(N+1)}2\right).
\]
After fixing $N$ modulo $2b$, the phase $(-1)^N\omega^{N(N+1)/2}$ is fixed.  Combining this identity with \eqref{eq:psi-cusp-final} gives \eqref{eq:Lambda-final-asym} with
\[
        f_b(\tau)=\frac{\pi^2}{2b^2\tau}-\frac{\tau}{2},
\]
and all remaining factors absorbed into a holomorphic non-vanishing $A_N$ satisfying \eqref{eq:AN-bound}.
\end{proof}

\end{document}
